%%%%%%%%%%%%%%%%%%%%%%%%%%%%%%%%%%%%%%%%%
% C++ Cheat Sheet Edition (3-Column Poster Layout)
%%%%%%%%%%%%%%%%%%%%%%%%%%%%%%%%%%%%%%%%%

\documentclass[final]{beamer}

\usepackage[scale=1.20]{beamerposter}
\usetheme{confposter}

\usepackage{graphicx}
\usepackage{booktabs}
\usepackage{hyperref}
\usepackage{amsmath}


\setbeamercolor{block title}{fg=ngreen,bg=white}
\setbeamercolor{block body}{fg=black,bg=white}

\setlength{\paperwidth}{36in}
\setlength{\paperheight}{48in}
\setlength{\topmargin}{-0.5in}


\usepackage{listings}
\usepackage{xcolor}

% Define a pale gray background
\definecolor{lightgray}{gray}{0.95}

\lstset{
    basicstyle=\ttfamily\footnotesize,   % smaller monospaced font
    breaklines=true,
    columns=fullflexible,
   % frame=single,                        % box around code
   % backgroundcolor=\color{lightgray},   % pale gray background
   % rulecolor=\color{black},             % border color
   % tabsize=4,
   % showspaces=false,
   % showstringspaces=false
}


\title{C++ QUICK REFERENCE / C++ CHEATSHEET}
\author{Based on Phillip M. Duxbury's C++ Cheatsheet and Morten Nobel-Jørgensen edits}

\begin{document}

\begin{frame}[t,fragile]

\begin{columns}[t,totalwidth=\paperwidth, colsep=0.015\paperwidth]

% ---------------- Column 1 ----------------
\begin{column}{0.32\paperwidth}

\begin{block}{Preprocessor}
\begin{lstlisting}[language=C++]
// Comment to end of line
/* Multi-line comment */
#include <stdio.h>         // Insert standard header file
#include "myfile.h"         // Insert file in current directory
#define X some text         // Replace X with some text
#define F(a,b) a+b          // Replace F(1,2) with 1+2
#define X \
 some text                 // Multiline definition
#undef X                    // Remove definition
#if defined(X)              // Conditional compilation (#ifdef X)
#else                       // Optional (#ifndef X or #if !defined(X))
#endif                      // Required after #if, #ifdef
\end{lstlisting}
\end{block}

\begin{block}{Literals}
\begin{lstlisting}[language=C++]
255, 0377, 0xff             // Integers (decimal, octal, hex)
2147483647L, 0x7fffffffl    // Long (32-bit) integers
123.0, 1.23e2               // double (real) numbers
'a', '\141', '\x61'         // Character (literal, octal, hex)
'\n', '\\', '\'', '\"'      // Newline, backslash, single quote, double quote
"string\n"                  // Array of characters ending with newline and \0
"hello" "world"             // Concatenated strings
true, false                 // bool constants 1 and 0
nullptr                     // Pointer type with the address of 0
\end{lstlisting}
\end{block}

\begin{block}{Declarations}
\begin{lstlisting}[language=C++]
int x;                      // Declare x to be an integer (value undefined)
int x=255;                  // Declare and initialize x to 255
short s; long l;            // Usually 16 or 32 bit integer (int may be either)
char c='a';                 // Usually 8 bit character
unsigned char u=255;
signed char s=-1;           // char might be either
unsigned long x = 0xffffffffL; // short, int, long are signed
float f; double d;          // Single or double precision real (never unsigned)
bool b=true;                // true or false, may also use int (1 or 0)
int a, b, c;                // Multiple declarations
int a[10];                  // Array of 10 ints (a[0] through a[9])
int a[]={0,1,2};            // Initialized array
int a[2][2]={ {1,2},{4,5} };  // Array of array of ints
char s[]="hello";           // String (6 elements including '\0')
std::string s = "Hello"     // Creates string object with value "Hello"
std::string s = R"(Hello\nWorld)"; // Creates string object with value "Hello\nWorld"
int* p;                     // p is a pointer to int
char* s="hello";            // s points to unnamed array containing "hello"
void* p=nullptr;            // Address of untyped memory
int& r=x;                   // r is a reference to (alias of) int x
enum weekend {SAT,SUN};     // weekend is a type with values SAT and SUN
enum weekend day;           // day is a variable of type weekend
enum weekend{SAT=0,SUN=1};  // Explicit representation as int
enum {SAT,SUN} day;         // Anonymous enum
enum class Color {Red,Blue};// Strict enum type
Color x = Color::Red;       // Assign Color x to red
typedef char* String;       // String s; means char* s;
const int c=3;              // Constants must be initialized
const int* p=a;             // Contents of p are constant
int* const p=a;             // p (not contents) is constant
const int* const p=a;       // Both p and contents constant
const int& cr=x;            // cr cannot assign to change x
int8_t,uint8_t,int16_t,uint16_t,int32_t,uint32_t,int64_t,uint64_t // Fixed length types
auto it = m.begin();        // Declares it to the result of m.begin()
auto const param = config["param"]; // Declares const result
auto& s = singleton::instance();   // Declares reference to result
\end{lstlisting}
\end{block}

\begin{block}{Storage Classes}
\begin{lstlisting}[language=C++]
int x;                      // Auto (memory exists only while in scope)
static int x;               // Global lifetime even if local scope
extern int x;               // Declared elsewhere
\end{lstlisting}
\end{block}

\begin{block}{Statements}
\begin{lstlisting}[language=C++]
x=y;                        // Every expression is a statement
int x;                      // Declarations are statements
;                           // Empty statement
{ int x; }                  // Block scope
if (x) a; else if (y) b; else c;
while (x) a;
for (x; y; z) a;
for (x : y) a;              // Range-based for
do a; while (x);
switch (x){ case X1:a; case X2:b; default:c; }
break; continue; return x;
try { a; } catch (T t) { b; } catch (...) { c; }
\end{lstlisting}
\end{block}

\begin{block}{Functions}
\begin{lstlisting}[language=C++]
int f(int x, int y);        // Function taking 2 ints returning int
void f();                   // Procedure taking no args
void f(int a=0);            // Default argument
inline f();                 // Suggest inline optimization
f() { statements; }         // Function definition
T operator+(T x, T y);      // Overload +
T operator-(T x);            // Unary -
T operator++(int);           // Postfix ++
extern "C" { void f(); }    // C linkage
int main() { statements; }     
int main(int argc, char* argv[]) { statements; }
\end{lstlisting}
\end{block}

\end{column}

% ---------------- Column 2 ----------------
\begin{column}{0.32\paperwidth}

\begin{block}{Expressions}
\begin{lstlisting}[language=C++]
T::X, N::X, ::X             // Name lookup
t.x, p->x, a[i], f(x,y), T(x,y)
x++, x--                    // Postfix increment/decrement
++x, --x                    // Prefix increment/decrement
typeid(x), typeid(T)
dynamic_cast<T>(x), static_cast<T>(x), reinterpret_cast<T>(x), const_cast<T>(x)
sizeof x, sizeof(T)
~x, !x, -x, +x, &x, *p
new T, new T(x,y), new T[x]
delete p, delete[] p
(T) x                       // C-style cast
x * y, x / y, x % y
x + y, x - y, x << y, x >> y
x < y, x <= y, x > y, x >= y
x & y, x ^ y, x | y
x && y, x || y
x = y, x += y, x -= y, x *= y, x /= y, x <<= y, x >>= y, x &= y, x |= y, x ^= y
x ? y : z
throw x
x , y                       // evaluates x then y
\end{lstlisting}
\end{block}

\begin{block}{Classes}
\begin{lstlisting}[language=C++]
class T {
private: int x;
protected: int y;
public:
    void f(); void g() { return; }
    void h() const;
    int operator+(int y);
    int operator-();
    T():x(1){}             // Constructor
    T(const T& t):x(t.x){} // Copy constructor
    T& operator=(const T& t){ x=t.x; return *this; } // Assignment
    ~T();                   // Destructor
    explicit T(int a);      // explicit constructor
    T(float x):T((int)x){}  // Delegate constructor
    operator int() const {return x;}
    friend void i();
    friend class U;
    static int y; 
    static void l();
    class Z {}; 
    typedef int V;
};
struct T{ virtual void i(); virtual void g()=0; }; 
class U: public T { void g(int) override; };
class V: private T {};
class W: public T, public U {};
class X: public virtual T {};
\end{lstlisting}
\end{block}

\begin{block}{Templates}
\begin{lstlisting}[language=C++]
template<class T> T f(T t);
template<class T> class X{ X(T t); };
template<class T> X<T>::X(T t) {}
X<int> x(3);
template<class T, class U=T, int n=0> class Y{};
\end{lstlisting}
\end{block}

\begin{block}{Namespaces}
\begin{lstlisting}[language=C++]
namespace N { class T {}; }
N::T t;
using namespace N;
\end{lstlisting}
\end{block}

\begin{block}{Memory Management}
\begin{lstlisting}[language=C++]
#include <memory>
shared_ptr<int> x;         // Empty shared_ptr
x = make_shared<int>(12);  // Allocate on heap
shared_ptr<int> y = x;     // Shared ownership
cout << *y;                // Dereference
y.reset();                 // Remove ownership
weak_ptr<int> w = y;       // Weak reference
if(shared_ptr<int> s = w.lock()){ cout << *s; }
unique_ptr<int> z; z = make_unique<int>(16);
unique_ptr<int> q; q = move(z);
if(z==nullptr) cout << "Z null";
cout << *q;
\end{lstlisting}
\end{block}

\begin{block}{cmath}
\begin{lstlisting}[language=C++]
#include <cmath>
sin(x); cos(x); tan(x);
asin(x); acos(x); atan(x);
atan2(y,x);
sinh(x); cosh(x); tanh(x);
exp(x); log(x); log10(x);
pow(x,y); sqrt(x);
ceil(x); floor(x);
fabs(x); fmod(x,y);
\end{lstlisting}


\end{block}

\end{column}

% ---------------- Column 3 ----------------
\begin{column}{0.32\paperwidth}

\begin{block}{cassert}
\begin{lstlisting}[language=C++]
#include <cassert>
assert(e);               // Abort if e is false
#define NDEBUG             // Disable asserts
\end{lstlisting}
\end{block}

\begin{block}{iostream}
\begin{lstlisting}[language=C++]
#include <iostream>
cin >> x >> y;           // Read input
cout << "x=" << 3 << endl; // Write output
cerr << x << y << flush;  // stderr
c = cin.get();            // getchar
cin.get(c);
cin.getline(s,n,'\n');
\end{lstlisting}
\end{block}

\begin{block}{fstream}
\begin{lstlisting}[language=C++]
#include <fstream>
ifstream f1("filename"); if(f1) f1 >> x;
f1.get(s); f1.getline(s,n);
ofstream f2("filename"); if(f2) f2 << x;
\end{lstlisting}
\end{block}

\begin{block}{string}
\begin{lstlisting}[language=C++]
#include <string>
string s1, s2="hello";
s1 += s2 + ' ' + "world";
s1 == "hello world";
s1[0]; s1.substr(m,n); s1.c_str();
s1 = to_string(12.05);
getline(cin,s);
\end{lstlisting}
\end{block}

\begin{block}{vector / deque}
\begin{lstlisting}[language=C++]
#include <vector>
vector<int> a(10); vector<int> b{1,2,3};
a.push_back(3); a.back()=4; a.pop_back();
a.front(); a.at(20)=1;  // safe vs unsafe access
for(int& p:a) p=0;      // C++11
for(vector<int>::iterator p=a.begin(); p!=a.end(); ++p) *p=0; // C++03
vector<int> b(a.begin(), a.end());
vector<T> c(n,x); T d[10]; vector<T> e(d,d+10);

#include <deque>
deque<int> d;
d.push_front(x); d.pop_front();
\end{lstlisting}
\end{block}

\begin{block}{pair / map / set}
\begin{lstlisting}[language=C++]
#include <utility>
pair<string,int> a("hello",3);
a.first; a.second;

#include <map>
map<string,int> a; a["hello"]=3;
for(auto &p:a) cout << p.first << p.second;

#include <unordered_map>
unordered_map<string,int> a; a["hello"]=3;

#include <set>
set<int> s; s.insert(123);
if(s.find(123) != s.end()) s.erase(123);

#include <unordered_set>
unordered_set<int> s; s.insert(123);
if(s.find(123) != s.end()) s.erase(123);
\end{lstlisting}
\end{block}

\end{column}

\end{columns}

\end{frame}

\end{document}
